% !TEX program = xelatex
\documentclass[10pt, a4paper,twocolumn]{article}
% ===== Essential Packages =====
\usepackage[utf8]{inputenc}       % UTF-8 encoding
\usepackage[T1]{fontenc}          % Better font encoding
\usepackage{lmodern}              % Modern Latin Modern font
\usepackage{ctex}                 % Chinese support (must be loaded after fontenc)
\usepackage{subcaption}
% ===== Math & Symbols =====
\usepackage{amsmath}              % Advanced math
\usepackage{amssymb}              % Additional math symbols
\usepackage{mathrsfs}             % Ralph Smith's Formal Script
\usepackage{braket}               % Dirac bra-ket notation

% ===== Graphics & Tables =====
\usepackage{graphicx}             % Image support
\usepackage{float}                % Better float control
\usepackage{wrapfig}              % Wrapped figures
\usepackage{tikz}                 % Vector graphics
\usepackage[table]{xcolor}        % Colored tables
\usepackage{array}                % Enhanced tables
\usepackage{multirow}             % Multirow tables
\usepackage{tabularx}             % Adjustable-width tables
\usepackage{booktabs}             % Professional table lines
\usepackage{siunitx}              % SI units in tables
\usepackage{caption}              % Better captions
\usepackage{adjustbox}            % Box resizing

% ===== Code Listings =====
\usepackage{listings}             % Code blocks
\usepackage{tcolorbox}            % Colored boxes
\tcbuselibrary{listings, breakable} % Listings in tcolorbox

% ===== Misc =====
\usepackage{lipsum}               % Dummy text
\usepackage{footmisc}             % Better footnotes
\usepackage[margin=1in]{geometry} % Page margins
\usepackage{diagbox}
\usepackage{ctex}
\author{张世航}
\title{磁光克尔效应}
\begin{document}
	\maketitle
	\section{磁滞回线}
	\begin{figure}[H]
		\centering
		\includegraphics[width=0.7\linewidth]{1}
		\caption{磁滞回线}
		\label{fig:1}
	\end{figure}
	\section{角度定标}
    \begin{figure}[H]
        \centering
		\includegraphics[width=0.7\linewidth]{2}
		\caption{角度定标}
		\label{fig:2}
    \end{figure}
    \section{螺旋测微器定标}
    \begin{figure}[H]
        \centering
		\includegraphics[width=0.7\linewidth]{3}
		\caption{螺旋测微器定标}
		\label{fig:3}
    \end{figure}
    \section{磁感应强度定标}
    \begin{figure}[H]
        \centering
		\includegraphics[width=0.7\linewidth]{4}
		\caption{磁感应强度定标}
		\label{fig:4}
    \end{figure}
	\section{数据处理}
	我们注意到磁滞回线的横坐标为\(U_B(mV)\),纵坐标为S.E.V采样值\(U_{S.E.V}/mV\).\\
	我们得到角度的标定最小二乘法拟合为
	\begin{equation}
		L(mm) = 16.60661\theta -24.6661
	\end{equation}
	我们得到螺旋测微器定标最小二乘法拟合为
	\begin{equation}
		U_{S.E.V}(V)= 6.09475L(mm)  -1.7143
	\end{equation}
	我们得到磁感应强度定标最小二乘法拟合为
	\begin{equation}
		B(mT)= 1250.00000U_B(V)  -1273.5000
	\end{equation}
	则
	\begin{equation}
		U_{S.E.V}(V) =6.09475(16.60661\theta -24.6661)-1.7143 
	\end{equation}
	可以反解出\(\theta\)的表达式
	\begin{equation}
		\theta = 1.50226 + 0.00988015 U_{S.E.V}(V)
	\end{equation}	
	画出此时的磁滞回线
	\begin{figure}[H]
        \centering
		\includegraphics[width=0.7\linewidth]{1.1}
		\caption{修正的磁滞回线}
		\label{fig:1.1}
    \end{figure}
	我们可以计算角度的变化为0.0009089738000001013,换成角度为
	\begin{equation}
		\Delta\theta = \frac{0.0009089738000001013\times 180^o}{\pi}= 0.05208036331636961^o
	\end{equation}
	根据图片我们可以看到矫顽力值为-10.545mT.
\end{document} 